\begin{textsample}{POS Dim 3 – rr – Score 19.00 – rr/rr\_inf\_25.txt}  \label{ex:f3_pos_152}
3.14 TRANSIÇÃO DA EDUCAÇÃO \textbf{INFANTIL} PARA O ENSINO FUNDAMENTAL A Educação \textbf{Infantil} é a primeira etapa da Educação Básica que educa e cuida de \textbf{crianças} de 0 a 5 anos de idade e passa por diferentes momentos de transição.
Nesse contexto, o primeiro processo de transição acontece de casa para a instituição de Educação \textbf{Infantil} e, posteriormente, a transição da \textbf{creche} para a \textbf{pré-escola} e, posteriormente, da \textbf{pré-escola} ao Ensino Fundamental.
Diante disso, é necessário que a escola busque a participação da família desde cedo para que aconteça, nesses processos de transição, o desenvolvimento integral nos aspectos éticos, políticos e estéticos, como tratam as Diretrizes Curriculares da Educação \textbf{Infantil} na Resolução nº 5, de dezembro de 2009.
A transição entre essas duas etapas da Educação Básica requer muita atenção, para que haja \textbf{equilíbrio} entre as mudanças introduzidas, garantindo integração e continuidade dos processos de aprendizagens das \textbf{crianças}, respeitando suas singularidades e as diferentes relações que elas estabelecem com os conhecimentos, assim como a natureza das mediações de cada etapa.
Torna -se necessário estabelecer estratégias de acolhimento e adaptação tanto para as \textbf{crianças} quanto para os docentes, de modo que a nova etapa se construa com base no que a \textbf{criança} sabe e é capaz de fazer, em uma perspectiva de continuidade de seu percurso educativo.
Para isso, as informações contidas em relatórios, portfólios ou outros \textbf{registros} que evidenciem os processos vivenciados pelas \textbf{crianças} ao longo de sua trajetória na Educação \textbf{Infantil} podem contribuir para a compreensão da história de vida escolar de cada aluno do Ensino Fundamental.
Conversas ou visitas e troca de materiais entre os professores das escolas de Educação \textbf{Infantil} e de Ensino Fundamental – Anos Iniciais também são importantes para facilitar a inserção das \textbf{crianças} nessa nova etapa da vida escolar.
Além disso, para que as \textbf{crianças} superem com sucesso os desafios da transição, é indispensável um \textbf{equilíbrio} entre as mudanças introduzidas, a continuidade das aprendizagens e o acolhimento afetivo, de modo que a nova etapa se construa com base no que os educandos sabem e são capazes de fazer, evitando a fragmentação e a descontinuidade do trabalho pedagógico.
Nessa direção, considerando os direitos e os objetivos de aprendizagem e desenvolvimento, apresenta -se a síntese das aprendizagens esperadas em cada campo de experiências.
Essa síntese deve ser compreendida como elemento balizador e indicativo de objetivos a ser explorados em todo o segmento da Educação \textbf{Infantil}, e que serão ampliados e aprofundados no Ensino Fundamental, e não como condição ou pré-requisito para o acesso ao Ensino Fundamental.
SÍNTESE DE APRENDIZAGENS O EU, O OUTRO E O NÓS - Respeitar e expressar sentimentos e emoções. - Atuar em grupo e \textbf{demonstrar} interesse em construir novas relações, respeitando a diversidade e solidarizando -se com os outros. - Conhecer e respeitar regras de convívio social, manifestando respeito pelo outro.
CORPO, \textbf{GESTOS} E MOVIMENTOS - Reconhecer a importância de ações e situações do cotidiano que contribuem para o \textbf{cuidado} de sua saúde e a manutenção de ambientes saudáveis. - Apresentar autonomia nas práticas de \textbf{higiene}, alimentação, vestir -se e no \textbf{cuidado} com seu bem-estar, valorizando o próprio corpo. - Utilizar o corpo intencionalmente ( com criatividade, controle e adequação ) como instrumento de interação com o outro e com o meio. - Coordenar suas habilidades manuais.
TRAÇOS, \textbf{SONS} E FORMAS - Discriminar os diferentes tipos de \textbf{sons} e ritmos e interagir com a música, percebendo -a como forma de expressão individual e coletiva. - Expressar -se por meio das artes visuais, utilizando diferentes materiais. - Relacionar -se com o outro empregando \textbf{gestos}, palavras, \textbf{brincadeiras}, jogos, \textbf{imitações}, observações e expressão corporal. \textbf{ESCUTA}, FALA, PENSAMENTO E IMAGINAÇÃO - Expressar ideias, \textbf{desejos} e sentimentos em distintas situações de interação, por diferentes meios. - Argumentar e relatar fatos oralmente, em sequência temporal e causal, organizando e adequando sua fala ao contexto em que é produzida. - Ouvir, compreender, \textbf{contar}, recontar e criar narrativas. - Conhecer diferentes gêneros e portadores textuais, \textbf{demonstrando} compreensão da função social da escrita e reconhecendo a leitura como fonte de prazer e informação.
ESPAÇOS, TEMPOS, QUANTIDADES, RELAÇÕES E TRANSFORMAÇÕES - Identificar, nomear adequadamente e comparar as propriedades dos objetos, estabelecendo relações entre eles. - Interagir com o meio ambiente e com fenômenos naturais ou artificiais, \textbf{demonstrando} \textbf{curiosidade} e \textbf{cuidado} com relação a eles. - Utilizar vocabulário relativo às noções de grandeza ( maior, \textbf{menor}, igual etc. ), espaço ( dentro e fora ) e medidas ( comprido, curto, grosso, fino ) como meio de comunicação de suas experiências. - Utilizar unidades de medida ( dia e noite : dias, semanas, \textbf{meses} e ano ) e noções de tempo ( presente, passado e futuro, antes, agora e depois ), para responder a necessidades e questões do cotidiano. - Identificar e registrar quantidades por meio de diferentes formas de representação ( contagens, desenhos, símbolos, escritas de números, organização de gráficos básicos etc. ).

% matched lemmas: brincadeira, contar, creche, criança, cuidado, curiosidade, demonstrar, desejo, equilíbrio, escuta, gesto, higiene, imitação, infantil, mês, pequeno, pré-escola, registro, som
\end{textsample}
